\fchapter{Ej 4 Pág 74}

\fsection{\boldit{\textcolor{waveBlue2}{Analisis.}}{ Teniendo en cuenta cómo el uso de la nube puede afectar a la rentabilidad de la empresa, te proponemos varios ejemplos para que los relaciones con los diferentes aspectos clave trabajados en este punto:}}

\begin{solution}
	\boldit{Modelo de costes eficiente / Flexibilidad de pago por uso (PAYG):}
	\begin{itemize}
		\item Una aplicación móvil de streaming de vídeo paga a su proveedor de nube solo por el ancho de banda y almacenamiento que consume, adaptando sus gastos a los ingresos generados.
		\item Un fabricante ajusta su uso de la nube trimestralmente, basándose en la demanda del mercado, evitando contratos a largo plazo.
	\end{itemize}

	\boldit{Reducción de costes operativos / Ahorros en mantenimiento y actualizaciones:}
	\begin{itemize}
		\item Una empresa minorista traslada su infraestructura de TI a la nube, disminuyendo gastos en mantenimiento de hardware y administración de sistemas, lo que le permite enfocarse más en mejorar la experiencia del cliente.

		\item Una agencia de marketing digital utiliza software en la nube que se actualiza automáticamente.
	\end{itemize}
	\boldit{Escalabilidad y adaptabilidad:}
	\begin{itemize}
		\item Un comercio electrónico ajusta sus recursos en la nube durante eventos de venta masiva, como el Black Friday, para manejar picos de demanda sin incurrir en costes por capacidad no utilizada.
	\end{itemize}

	\boldit{Mejora de la rentabilidad a largo plazo / Reducción de riesgos financieros}
	\begin{itemize}
		\item Un banco lanza una nueva aplicación financiera en semanas utilizando la nube, acelerando el tiempo de comercialización y obteniendo ingresos más rápidamente.
	\end{itemize}

	\boldit{Optimización de recursos y tiempos de implementación}
	\begin{itemize}
		\item Una startup tecnológica usa almacenamiento y computación en la nube, pagando solo por el espacio y la potencia que usa, evitando así la compra de servidores propios.
		\item Una pequeña editorial utiliza herramientas avanzadas de análisis de datos en la nube para entender mejor a su audiencia.
	\end{itemize}

	\boldit{Focalización en la innovación y el core business}
	\begin{itemize}
		\item Una empresa de biotecnología externaliza su infraestructura de TI a la nube, permitiéndole concentrarse en la investigación y desarrollo de nuevos medicamentos.
		\item Al adoptar servicios en la nube, una firma de consultoría logra reducir costes y acceder a tecnología avanzada.
	\end{itemize}

\end{solution}
